\bookchapter{第十八章：数据备份与迁移}

\chapterintro{我一个邻居，手机掉进水池里捞出来已经开不了机了。她最心疼的不是手机本身，而是里面三年没备份过的照片——孩子从出生到会走路的全部记录。听她说完我心里特别不是滋味，因为这件事完全是可以避免的。这章我们就认认真真地给数据上一道保险。}

\booksection{一、先搞清楚：你手机里哪些东西丢了就没了}

不同数据的“生存状态”不一样，你得分清楚：

\begin{itemize}
  \item \textbf{存在手机本地、丢了就真的没了的}：照片视频（没开云同步的话）、备忘录、短信、通话记录、下载的文件。
  \item \textbf{存在App服务器上、换手机重新登录就有}：微信好友列表、淘宝订单、银行App里的交易记录。这些东西存储在对应公司的服务器上，不是存在你手机里的。
  \item \textbf{处于灰色地带、需要特别操作的}：微信聊天记录（虽然微信服务器存了，但你自己没法从服务器下载）、手机银行绑定的设备凭证、两步验证App里的密钥（Google Authenticator那种）。
\end{itemize}

所以你要重点保护的，是第一类和第三类。

\booksection{二、整机备份：一次搞定所有东西}

整机备份就像给手机拍一张完整的“快照”——当前系统设置、桌面布局、App数据、照片、通讯录，全部打包。

\booksubsection{方案一：电脑本地备份（最可靠，推荐有电脑的人）}

\begin{itemize}
  \item \textbf{iPhone}：把手机用数据线连到电脑上，打开iTunes（Windows旧版）或Finder（Mac），在设备页面里选"立即备份"。注意要勾上"给备份加密"，否则密码和健康数据不会被备份进去。
  \item \textbf{安卓}：各品牌有自己的电脑端软件，比如华为手机助手、小米助手、三星Smart Switch。连上电脑按提示操作就行。
  \item 电脑备份的好处是不用花钱买云空间，而且恢复速度比云端恢复快得多。缺点是需要你记得定期去插线备份——大部分人做不到这一点。
\end{itemize}

\booksubsection{方案二：云备份（最省心，推荐所有人）}

\begin{itemize}
  \item \textbf{iPhone}：设置 你的名字 iCloud iCloud云备份 打开。手机插着电源、连Wi-Fi、锁屏的时候会自动备份。免费的5GB大概率不够，建议花6块钱升级到50GB。
  \item \textbf{安卓}：各家品牌云都有备份功能。华为手机去设置搜"云备份"，小米去"小米云服务"，OPPO/vivo也类似。另外还有一个谷歌自带的备份（设置 系统 备份），这个也打开它。
  \item 云备份最大的好处是"不用你操心"，自动搞定。最大的坏处是恢复的时候要下载很久，而且免费空间一般不够用。
\end{itemize}

\booksection{三、分类备份策略：重要的单独另存一份}

整机备份是兜底的，但有些东西太重要了，值得你单独另存一份：

\booksubsection{照片和视频（最珍贵的数字资产）}

\begin{itemize}
  \item \textbf{双保险方案}：开iCloud照片/Google相册自动同步，同时每隔几个月手动把照片复制一份到电脑或移动硬盘上。这样云端一份、本地一份，互为备份。
  \item \textbf{不想花钱的方案}：Google相册免费给15GB，可以设置为"压缩画质"来省空间。或者每隔一个月用数据线把照片全部拷到电脑上。
\end{itemize}

\booksubsection{通讯录（容易忽略但极其重要）}

\begin{itemize}
  \item iPhone通讯录默认和iCloud同步，只要开了同步就不用额外备份。但如果你想再安全一点，可以去iCloud.com导出vCard文件存电脑上。
  \item 安卓通讯录通常和谷歌或品牌云同步。去通讯录App的设置里导出.vcf文件，就是一份完整的备份。
\end{itemize}

\booksubsection{微信聊天记录（单独处理，重中之重）}

这个必须单独拿出来说，因为整机备份里微信聊天记录很多时候不完整，而且换机工具根本搬不动它。

\begin{itemize}
  \item \textbf{方法一：迁移到新手机}。微信 我 设置 通用 聊天记录迁移与备份 迁移 迁移到手机/平板。新旧手机连同一个Wi-Fi，扫码就能传。注意：迁移完成之前，旧手机别退出微信。
  \item \textbf{方法二：备份到电脑}。电脑上装微信PC版，左下角三条杠 迁移与备份 备份与恢复 备份聊天记录到电脑。之后换手机了可以再恢复到新手机上。
  \item \textbf{方法三：备份到云端}。微信也支持备份到iCloud或手机云空间，但注意：这个备份只能在你当前手机上恢复，不能跨设备还原。换手机的话，还是得用前两种方法。
\end{itemize}

\booksection{四、换机指南：从旧手机到新手机}

换新手机是高兴事，但如果数据没搬过去就变成头疼事了。三种情况分别对待：

\booksubsection{同品牌升级（最简单）}

\begin{itemize}
  \item \textbf{iPhone 新iPhone}：开机激活时把旧手机放在旁边，会自动弹出"快速开始"，用旧手机扫新手机上的图案就行。所有数据、设置、App布局都会原封不动搬过去，包括Wi-Fi密码都会自动带过去。大概需要二三十分钟。
  \item \textbf{安卓同品牌}：用品牌自带的换机工具。华为用"手机克隆"，小米用"小米换机"，OPPO用"手机搬家"，vivo用"互传"。步骤都差不多：新旧手机装同一个App、扫码连接、勾选要搬的东西、开传。
\end{itemize}

\booksubsection{跨品牌换机（安卓换安卓，不同品牌）}

用第三方换机工具或者用旧品牌自带的换机App（大部分品牌的换机App也支持换到其他安卓手机）。能搬走的有：通讯录、照片、短信、通话记录、部分App。搬不走的：App内的登录状态（你得重新登录）、微信聊天记录（用微信自己的迁移功能）。

\booksubsection{跨平台换机（安卓换iPhone，或反过来）}

\begin{itemize}
  \item \textbf{安卓 iPhone}：苹果有个"转移到iOS"App，安卓手机上下载后，在新iPhone激活时选择"从安卓设备转移数据"。能搬走通讯录、照片、短信等基础数据，但App数据搬不了。微信聊天记录还得用微信自己的迁移功能。
  \item \textbf{iPhone 安卓}：各安卓品牌有对应的迁移工具，但说实话，跨平台迁移体验不如同平台好。最重要的事先单独做了：通讯录导出、照片备份、微信聊天记录迁移。剩下的慢慢来。
\end{itemize}

\booksection{五、备份的黄金法则：三二一原则}

这是数据安全领域的一个经典原则，简单易记：

\begin{itemize}
  \item \textbf{3}：重要数据至少保留三份（一份原件，两份备份）。
  \item \textbf{2}：备份用两种不同的存储介质（比如云端一份，移动硬盘一份）。
  \item \textbf{1}：至少有一份备份放在不同的物理位置（防止火灾、盗窃等同时毁掉所有数据）。
\end{itemize}

对普通人来说，做到“云端一份、电脑一份”就已经算很好了。不用搞太复杂，关键是你真的去做了。

\begin{tcbwarn}{微信聊天记录特别注意}
微信聊天记录无法通过任何手机换机工具迁移！它不在整机备份的覆盖范围内。必须用微信自带的“聊天记录迁移”功能，或备份到电脑版微信。很多人在换完新手机之后才发现聊天记录全没了，那时候已经晚了——旧手机可能已经恢复出厂设置了。

\end{tcbwarn}

\begin{tcbpractice}{练一练：做一次完整的备份演练}
\begin{enumerate}
  \item 检查iCloud云备份或品牌云备份是否开启，确认上一次备份时间。
  \item 手动备份照片：用数据线连接电脑，把DCIM和Pictures文件夹复制到电脑。
  \item 备份通讯录：iPhone用户去iCloud.com导出vCard；安卓用户在联系人App设置里导出.vcf文件。
  \item 备份微信聊天记录：打开微信 我 设置 通用 聊天记录迁移与备份 备份到电脑。
  \item 把导出的通讯录vcf和照片文件夹复制一份到移动硬盘或U盘，完成"三二一"里的"一份离线在别处"。
\end{enumerate}

\end{tcbpractice}
